\documentclass[aspectratio=169,10pt,spanish]{beamer}

\usepackage[utf8]{inputenc}
\usepackage[T1]{fontenc}
\usepackage{lmodern}
\usepackage{helvet}
\renewcommand{\familydefault}{\sfdefault}
\usepackage[spanish,es-tabla,es-nodecimaldot]{babel}
\usepackage{tikz}
\usetikzlibrary{arrows.meta,positioning,calc}
\usepackage[most]{tcolorbox}
\usepackage{amsmath,amssymb}
\usepackage{booktabs,tabularx,array}
\usepackage{ragged2e}

\setbeamertemplate{navigation symbols}{}
\setbeamercolor{background canvas}{bg=white}
\definecolor{bpNavy}{RGB}{28,43,54}
\definecolor{bpBlue}{RGB}{43,91,132}
\definecolor{bpGold}{RGB}{200,157,66}
\definecolor{bpRust}{RGB}{179,92,62}
\definecolor{bpDark}{RGB}{33,37,41}
\definecolor{bpLightGray}{RGB}{245,247,250}
\definecolor{bpGridLine}{RGB}{225,230,238}
\definecolor{bpGreen}{RGB}{46,125,50}
\definecolor{bpRed}{RGB}{170,56,49}

\setbeamercolor{block title}{fg=white,bg=bpBlue}
\setbeamercolor{block body}{fg=bpDark,bg=bpLightGray}
\setbeamercolor{alerted text}{fg=bpRust}
\setbeamertemplate{itemize item}{\color{bpGold}\small\raise0.2ex\hbox{\textbullet}}
\setbeamertemplate{itemize subitem}{\color{bpBlue}\tiny\raise0.2ex\hbox{\textbullet}}

\setbeamertemplate{frametitle}{
  \begin{tikzpicture}[remember picture,overlay]
    \fill[bpNavy] (current page.north west) rectangle ([yshift=-1.05cm]current page.north east);
    \node[anchor=west,text=white,font=\bfseries\sffamily\normalsize] at ([xshift=0.6cm,yshift=-0.52cm]current page.north west) {\insertframetitle};
    \node[anchor=east,text=white!75,font=\tiny\sffamily] at ([xshift=-0.6cm,yshift=-0.52cm]current page.north east) {SEMINARIO DE INVESTIGACIÓN II};
  \end{tikzpicture}
  \vspace{0.24cm}
}
\setbeamertemplate{footline}{
  \begin{tikzpicture}[remember picture,overlay]
    \draw[bpGridLine,line width=0.8pt] ([yshift=0.5cm]current page.south west)--([yshift=0.5cm]current page.south east);
    \node[anchor=west,text=bpNavy!80!black,font=\tiny\sffamily] at ([xshift=0.6cm,yshift=0.25cm]current page.south west) {UNSAAC -- Ingeniería Informática y de Sistemas $\mid$ Br. Edmil Jampier Saire Bustamante};
    \node[anchor=east,text=bpNavy!80!black,font=\tiny\sffamily] at ([xshift=-0.6cm,yshift=0.25cm]current page.south east) {\insertframenumber{} / \inserttotalframenumber};
  \end{tikzpicture}
}
\newcommand{\fuente}[1]{\vfill{\tiny\color{bpNavy!75}\textbf{Base documental:} #1}}
\newcommand{\dato}[2]{\textbf{\color{bpBlue}#1} #2}
\newcommand{\etiqueta}[1]{\tcbox[colback=bpNavy,colframe=bpNavy,boxrule=0pt,arc=1mm,left=2mm,right=2mm,top=0.7mm,bottom=0.7mm]{\color{white}\scriptsize\bfseries #1}}

\begin{document}

% 1. Portada
\begin{frame}[plain]
  \begin{tikzpicture}[remember picture,overlay]
    \draw[step=0.5cm,bpGridLine,line width=0.25pt] (current page.south west) grid (current page.north east);
    \draw[bpNavy,line width=1.4pt] ([xshift=0.5cm,yshift=0.5cm]current page.south west) rectangle ([xshift=-0.5cm,yshift=-0.5cm]current page.north east);
    \fill[bpNavy] ([xshift=0.5cm,yshift=-1.55cm]current page.north west) rectangle ([xshift=-0.5cm,yshift=-0.5cm]current page.north east);
    \node[anchor=center,text=white,font=\large\bfseries\sffamily,align=center,text width=12.8cm] at ([yshift=-1.02cm]current page.north) {DEL CONTEXTO A LA EVALUACIÓN\\EL PROCESO COMPLETO DEL PLAN DE TESIS};
    \node[anchor=center,align=center,text width=12.8cm] at ([yshift=0.45cm]current page.center) {
      {\color{bpNavy}\Large\bfseries Diseño e implementación de una arquitectura SLM--RAG con QLoRA}\\[0.22cm]
      {\color{bpDark}\normalsize para la generación contextualizada en Quechua Collao}\\[0.45cm]
      \begin{tcolorbox}[width=11.8cm,colback=white,colframe=bpGold,arc=1.5mm,boxrule=1pt]
        \centering\small\textbf{Ruta de investigación aplicada:} contexto $\rightarrow$ evidencia $\rightarrow$ problema $\rightarrow$ objetivos $\rightarrow$ método $\rightarrow$ evaluación
      \end{tcolorbox}
    };
    \node[anchor=south,align=center,font=\scriptsize\sffamily,text=bpDark] at ([yshift=0.75cm]current page.south) {
      \textbf{Tesista:} Br. Edmil Jampier Saire Bustamante \quad $\mid$ \quad \textbf{Cátedra:} Seminario de Investigación II\\
      UNSAAC -- Cusco, 2026
    };
  \end{tikzpicture}
\end{frame}

% 2. Mapa del proceso
\begin{frame}{La ruta completa aplicada al tema}
\small
\begin{columns}[T,onlytextwidth]
\column{0.32\textwidth}
\begin{tcolorbox}[colback=bpLightGray,colframe=bpRust,title={1. Contexto}]
Quechua Collao, recursos limitados y meta de uso offline.
\end{tcolorbox}
\column{0.32\textwidth}
\begin{tcolorbox}[colback=bpLightGray,colframe=bpBlue,title={2. Evidencia}]
Corpus, diccionarios y estudios sobre recuperación y SLM.
\end{tcolorbox}
\column{0.32\textwidth}
\begin{tcolorbox}[colback=bpLightGray,colframe=bpRust,title={3. Problema}]
¿El SLM aprovecha el contexto recuperado para generar en Collao?
\end{tcolorbox}
\end{columns}
\vspace{0.28cm}
\begin{columns}[T,onlytextwidth]
\column{0.32\textwidth}
\begin{tcolorbox}[colback=bpLightGray,colframe=bpBlue,title={4. Objetivos}]
Curar datos, adaptar el modelo y comparar configuraciones.
\end{tcolorbox}
\column{0.32\textwidth}
\begin{tcolorbox}[colback=bpLightGray,colframe=bpGold,title={5. Justificación}]
Aporte lingüístico, técnico y práctico para ejecución local.
\end{tcolorbox}
\column{0.32\textwidth}
\begin{tcolorbox}[colback=bpLightGray,colframe=bpGold,title={6. Método}]
DSR, experimento E1/E2/E3 y métricas de calidad y recursos.
\end{tcolorbox}
\end{columns}
\vspace{0.18cm}
\centering\scriptsize Cada etapa debe dejar una decisión rastreable en el plan.
\end{frame}
% 3. Contexto
\begin{frame}{1 | Contexto: necesidad y oportunidad}
\begin{columns}[T,onlytextwidth]
\column{0.57\textwidth}
\begin{block}{Situación de investigación}
\begin{itemize}\small
  \item El plan se enfoca en Quechua Collao, una variedad que requiere recursos identificados y validados.
  \item El acceso a conocimiento confiable depende de corpus, diccionarios y documentos digitalizados.
  \item Un servicio exclusivamente en la nube no satisface el objetivo de inferencia offline.
  \item Un SLM compacto reduce requisitos de despliegue, pero su calidad en la variedad debe medirse.
\end{itemize}
\end{block}
\column{0.40\textwidth}
\begin{tcolorbox}[colback=bpNavy,colframe=bpNavy,arc=2mm]
\centering\color{white}\large\bfseries
Necesidad\\[2mm]
\color{white}\normalsize acceso local a respuestas contextualizadas\\[4mm]
\color{bpGold}\Large$\Downarrow$\\[2mm]
\color{white}\normalsize evidencia documental + generación lingüística
\end{tcolorbox}
\end{columns}
\fuente{Plan de tesis, introducción, problema y factibilidad.}
\end{frame}

% 4. Definición operativa del tema
\begin{frame}{El tema expresado como proceso + tecnología + contexto}
\centering
\begin{columns}[T,onlytextwidth]
\column{0.31\textwidth}
\begin{tcolorbox}[colback=bpLightGray,colframe=bpRust,title={Proceso}]
\centering\Large\color{bpRust}\textbf{P}\\[2mm]
\normalsize Diseñar e implementar una arquitectura y evaluar su generación contextualizada.
\end{tcolorbox}
\column{0.34\textwidth}
\begin{tcolorbox}[colback=bpLightGray,colframe=bpBlue,title={Tecnología}]
\centering\Large\color{bpBlue}\textbf{T}\\[2mm]
\normalsize SLM cuantizado + recuperación local (RAG) + adaptación eficiente (QLoRA).
\end{tcolorbox}
\column{0.31\textwidth}
\begin{tcolorbox}[colback=bpLightGray,colframe=bpGold,title={Contexto}]
\centering\Large\color{bpGold}\textbf{C}\\[2mm]
\normalsize Quechua Collao y ejecución en un entorno computacional offline.
\end{tcolorbox}
\end{columns}
\vspace{0.14cm}
\etiqueta{FORMULACIÓN DEL PLAN}
\vspace{0.04cm}
\begin{tcolorbox}[colback=white,colframe=bpNavy]
\centering\small Diseñar SLM--RAG con QLoRA para generar en Quechua Collao offline.
\end{tcolorbox}
\end{frame}

% 5. Marco conceptual
\begin{frame}{2 | Marco conceptual: qué hace cada componente}
\small
\begin{tabularx}{\textwidth}{>{\bfseries}p{2.45cm} X p{3.5cm}}
\toprule
Componente & Función prevista en la tesis & Pregunta de control\\
\midrule
SLM & Modelo compacto Qwen2.5-3B-Instruct cuantizado; genera la respuesta. & ¿Puede ejecutarse localmente con los recursos disponibles?\\
RAG local & Busca fragmentos documentales y los incorpora al prompt. & ¿El contexto recuperado es pertinente y trazable?\\
QLoRA & Ajusta adaptadores ligeros manteniendo congelado el modelo base cuantizado. & ¿Mejora el uso del contexto frente a RAG estándar?\\
Corpus bilingüe & Aporta evidencia textual en español y Quechua Collao. & ¿Se controla calidad, fuente y variedad dialectal?\\
\bottomrule
\end{tabularx}
\vspace{0.28cm}
\begin{tcolorbox}[colback=bpLightGray,colframe=bpGold]
\centering\small\textbf{Distinción clave:} RAG aporta conocimiento recuperable; QLoRA busca adaptar el comportamiento del modelo para utilizar el contexto.
\end{tcolorbox}
\end{frame}

% 6. Antecedentes: recursos
\begin{frame}{2 | Evidencia documentada: recursos para Quechua}
\begin{columns}[T,onlytextwidth]
\column{0.48\textwidth}
\begin{tcolorbox}[colback=bpLightGray,colframe=bpBlue,title={Corpus paralelo}]
\small Orcotoma Mormontoy et al. (2025) presentan un corpus paralelo Quechua--español construido a partir de publicaciones digitalizadas y OCR.\\[2mm]
\textbf{Aporte al plan:} semillas documentales y material para validar/curar.\\
\textbf{Cuidado:} OCR exige limpieza; no equivale a un corpus de entrenamiento listo.
\end{tcolorbox}
\column{0.48\textwidth}
\begin{tcolorbox}[colback=bpLightGray,colframe=bpGold,title={Recursos léxicos}]
\small Durán y Silberztein (2026) describen recursos léxicos y gramaticales para Quechua Sureño, incluidos verbos y formalización para procesamiento.\\[2mm]
\textbf{Aporte al plan:} apoyo para normalización y verificación léxica.\\
\textbf{Cuidado:} comprobar variedad, licencia y compatibilidad antes de integrar.
\end{tcolorbox}
\end{columns}
\fuente{Orcotoma Mormontoy et al. (2025); Durán y Silberztein (2026).}
\end{frame}

% 7. Antecedentes: RAG y SLM
\begin{frame}{2 | Evidencia documentada: recuperación y modelos compactos}
\small
\begin{columns}[T,onlytextwidth]
\column{0.32\textwidth}
\begin{block}{Court y Elsner (2024)}
Estudian traducción de Quechua Sureño al español con contexto. El efecto de proporcionar ejemplos depende del modelo, recuperación y contexto.
\end{block}
\column{0.32\textwidth}
\begin{block}{Hevia et al. (2025)}
Describen un tutor de IA con RAG y un modelo ligero para operar offline con pocos recursos; exploran estrategias de recuperación sobre libros universitarios.
\end{block}
\column{0.32\textwidth}
\begin{block}{Walusimbi et al. (2026)}
Proponen una arquitectura offline-first para entornos de conectividad limitada, con modelos locales cuantizados.
\end{block}
\end{columns}
\vspace{0.2cm}
\begin{tcolorbox}[colback=bpLightGray,colframe=bpNavy]
\centering\small\textbf{Lectura para esta tesis:} la recuperación de contexto y la inferencia offline son enfoques plausibles; su combinación no garantiza por sí sola fidelidad lingüística ni buen uso del contexto.
\end{tcolorbox}
\fuente{Estudios publicados; ver referencias al final.}
\end{frame}

% 8. QueEn
\begin{frame}{2 | Antecedente cercano: QueEn}
\begin{columns}[T,onlytextwidth]
\column{0.53\textwidth}
\begin{block}{Qué aporta}
\small Chen et al. (2024) presentan QueEn, un modelo de traducción Quechua--inglés que combina recuperación de información y ajuste eficiente de parámetros.
\begin{itemize}
  \item Antecedente de RAG + PEFT en una tarea que involucra Quechua.
  \item Da un punto de comparación conceptual para el enfoque híbrido.
  \item Su tarea y dirección de traducción no son idénticas a la generación contextualizada Collao planteada aquí.
\end{itemize}
\end{block}
\column{0.43\textwidth}
\begin{tcolorbox}[colback=bpLightGray,colframe=bpRust,title={Verificación bibliográfica}]
\small En la bibliografía local, la clave \texttt{Chen2024} apunta a un registro ACL distinto (NewsPaLM).\\[2mm]
El trabajo QueEn corresponde a arXiv:2412.05184.\\[2mm]
\textbf{Acción pendiente:} corregir la entrada BibTeX antes de la entrega formal del plan.
\end{tcolorbox}
\end{columns}
\fuente{Chen et al., \emph{QueEn: A Large Language Model for Quechua-English Translation} (2024), arXiv:2412.05184.}
\end{frame}

% 9. Síntesis/gap
\begin{frame}{De los antecedentes a la brecha investigable}
\centering
\begin{columns}[T,onlytextwidth]
\column{0.29\textwidth}
\begin{tcolorbox}[colback=bpLightGray,colframe=bpBlue,title={Lo que existe}]
Corpus y recursos léxicos; traducción con contexto; sistemas offline con SLM y RAG.
\end{tcolorbox}
\column{0.29\textwidth}
\begin{tcolorbox}[colback=bpLightGray,colframe=bpGold,title={Lo que falta comprobar}]
Cómo se comporta una integración concreta SLM--RAG--QLoRA para Quechua Collao, con evidencia local y restricciones de cómputo.
\end{tcolorbox}
\column{0.36\textwidth}
\begin{tcolorbox}[colback=bpNavy,colframe=bpNavy,title={Foco evaluable}]
\color{white}Comparar el uso de contexto recuperado entre modelo base, RAG estándar y RAG con QLoRA.
\end{tcolorbox}
\end{columns}
\vspace{0.24cm}
\begin{tcolorbox}[colback=white,colframe=bpRust]
\centering\small\textbf{Alcance de la afirmación:} es una brecha delimitada por este plan; no se afirma que nadie haya investigado previamente todos sus componentes.
\end{tcolorbox}
\end{frame}

% 10. Problema
\begin{frame}{3 | Problema de investigación}
\begin{block}{Problema general del plan}
\small Capacidad limitada de modelos compactos cuantizados para asimilar contextos recuperados en una arquitectura RAG offline orientada a generación en Quechua Collao.
\end{block}
\vspace{0.12cm}
\begin{columns}[T,onlytextwidth]
\column{0.32\textwidth}
\begin{tcolorbox}[colback=bpLightGray,colframe=bpRust,title={Datos y recuperación}]
Corpus heterogéneo; ruido OCR; necesidad de etiquetar fuentes y variedad lingüística.
\end{tcolorbox}
\column{0.32\textwidth}
\begin{tcolorbox}[colback=bpLightGray,colframe=bpBlue,title={Uso del contexto}]
El SLM puede recibir fragmentos pertinentes, pero no utilizarlos de forma completa o fiel.
\end{tcolorbox}
\column{0.32\textwidth}
\begin{tcolorbox}[colback=bpLightGray,colframe=bpGold,title={Restricción de cómputo}]
La solución debe operar localmente; adaptación y evaluación quedan sujetas a recursos acotados.
\end{tcolorbox}
\end{columns}
\vspace{0.16cm}
\centering\scriptsize El problema se traduce en una comparación experimental, no en una promesa de mejora.
\fuente{Plan de tesis, planteamiento del problema.}
\end{frame}

% 11. Preguntas
\begin{frame}{3 | Preguntas de investigación}
\small
\begin{tcolorbox}[colback=bpNavy,colframe=bpNavy]
\color{white}\textbf{Pregunta general}\\
¿De qué manera el diseño e implementación de una arquitectura SLM--RAG con QLoRA influye en la generación contextualizada en Quechua Collao dentro de un entorno computacional offline?
\end{tcolorbox}
\vspace{0.12cm}
\begin{tabularx}{\textwidth}{>{\bfseries}p{1.2cm} X}
\toprule
PE1 & ¿Cómo curar, estructurar e indexar una base documental bilingüe para recuperación local en Quechua Collao?\\
PE2 & ¿Cómo adaptar con QLoRA un SLM cuantizado para mejorar el uso del contexto recuperado dentro de las restricciones de hardware?\\
PE3 & ¿Cómo comparar fidelidad, calidad lingüística y eficiencia computacional frente al SLM base y al RAG estándar?\\
\bottomrule
\end{tabularx}
\fuente{Plan de tesis, preguntas general y específicas.}
\end{frame}

% 12. Objetivos
\begin{frame}{4 | Objetivo general y objetivos específicos}
\small
\begin{tcolorbox}[colback=bpNavy,colframe=bpNavy]
\color{white}\textbf{Objetivo general}\\
Diseñar e implementar una arquitectura SLM--RAG con QLoRA para la generación contextualizada en Quechua Collao dentro de un entorno computacional offline.
\end{tcolorbox}
\vspace{0.1cm}
\begin{columns}[T,onlytextwidth]
\column{0.32\textwidth}
\begin{block}{OE1 | Base documental}
Curar, estructurar e indexar una base de conocimiento bilingüe para recuperación local.
\end{block}
\column{0.32\textwidth}
\begin{block}{OE2 | Adaptación}
Implementar adaptadores QLoRA sobre un SLM cuantizado para asimilar contexto recuperado.
\end{block}
\column{0.32\textwidth}
\begin{block}{OE3 | Evaluación}
Comparar fidelidad, calidad lingüística y eficiencia frente a las configuraciones de referencia.
\end{block}
\end{columns}
\vspace{0.12cm}
\centering\scriptsize Los objetivos describen productos verificables: índice, adaptador y resultados comparativos.
\end{frame}

% 13. Hipótesis y variables
\begin{frame}{4 | Hipótesis, variables y sentido de la prueba}
\begin{columns}[T,onlytextwidth]
\column{0.50\textwidth}
\begin{tcolorbox}[colback=bpLightGray,colframe=bpGold,title={Hipótesis del plan}]
\small La arquitectura híbrida SLM--RAG con QLoRA mejorará la generación contextualizada respecto al SLM base y al RAG estándar, dentro de límites de cómputo offline.
\end{tcolorbox}
\begin{tcolorbox}[colback=white,colframe=bpRust,title={Estado}]
\small\textbf{Predicción por contrastar.} No es un resultado ni puede presentarse como hallazgo antes del experimento.
\end{tcolorbox}
\column{0.46\textwidth}
\begin{block}{Variable independiente}
Configuración del sistema:
\begin{itemize}\small
\item E1: SLM cuantizado, sin RAG.
\item E2: SLM cuantizado + RAG local.
\item E3: SLM cuantizado + RAG + QLoRA.
\end{itemize}
\end{block}
\begin{block}{Variables dependientes}
Fidelidad al contexto; calidad lingüística; eficiencia computacional.
\end{block}
\end{columns}
\fuente{Plan de tesis, hipótesis y diseño experimental.}
\end{frame}

% 14. Alineamiento
\begin{frame}{La matriz de consistencia en una sola vista}
\scriptsize
\begin{tabularx}{\textwidth}{p{2.15cm} p{3.55cm} p{3.35cm} X}
\toprule
\textbf{Problema específico} & \textbf{Objetivo asociado} & \textbf{Componente / producto} & \textbf{Evidencia de evaluación}\\
\midrule
PE1: evidencia documental para RAG & OE1: curar e indexar recursos bilingües & Corpus depurado, metadatos e índice local & Trazabilidad de fuentes; recuperación y precisión de contexto\\
\midrule
PE2: uso limitado del contexto & OE2: adaptar SLM con QLoRA & Adaptador entrenado y configuración E3 & Comparación E2 vs. E3 bajo el mismo banco de consultas\\
\midrule
PE3: impacto sin medir & OE3: comparar las tres configuraciones & Banco de evaluación y ejecución reproducible & Métricas lingüísticas, de fidelidad y de recursos\\
\bottomrule
\end{tabularx}
\vspace{0.25cm}
\begin{tcolorbox}[colback=bpLightGray,colframe=bpNavy]
\centering\small\textbf{Regla de consistencia:} cada pregunta tiene un objetivo; cada objetivo produce un artefacto o evidencia observable.
\end{tcolorbox}
\end{frame}

% 15. Justificación
\begin{frame}{5 | Justificación: por qué merece hacerse}
\begin{columns}[T,onlytextwidth]
\column{0.32\textwidth}
\begin{tcolorbox}[colback=bpLightGray,colframe=bpRust,title={Lingüística y social}]
Contribuye a explorar herramientas locales para una variedad que requiere recursos lingüísticos identificados y respetados.
\end{tcolorbox}
\column{0.32\textwidth}
\begin{tcolorbox}[colback=bpLightGray,colframe=bpBlue,title={Técnica}]
Evalúa si QLoRA puede mejorar el uso de contexto recuperado, además de medir el aporte propio de RAG.
\end{tcolorbox}
\column{0.32\textwidth}
\begin{tcolorbox}[colback=bpLightGray,colframe=bpGold,title={Práctica y académica}]
Deja un prototipo y un protocolo experimental comparativo que pueden reproducirse y ampliarse.
\end{tcolorbox}
\end{columns}
\vspace{0.28cm}
\begin{tcolorbox}[colback=white,colframe=bpNavy]
\centering\small\textbf{Delimitación ética:} fuentes públicas con atribución; control explícito de variante; resultados reportados sin prometer utilidad comunitaria no evaluada.
\end{tcolorbox}
\fuente{Plan de tesis, justificación y factibilidad.}
\end{frame}

% 16. Arquitectura
\begin{frame}{6 | Propuesta: flujo SLM--RAG con QLoRA}
\centering
\begin{tikzpicture}[
  node distance=0.38cm,
  caja/.style={draw=bpBlue,rounded corners=2pt,fill=bpLightGray,align=center,text width=2.35cm,minimum height=1.05cm,font=\scriptsize},
  arrow/.style={-{Stealth[length=2mm]},thick,draw=bpNavy}
]
\node[caja] (docs) {Fuentes bilingües\\y léxicos curados};
\node[caja,right=of docs] (chunk) {Limpieza,\\metadatos y chunks};
\node[caja,right=of chunk] (embed) {Embeddings\\e índice vectorial};
\node[caja,below=1.0cm of embed] (query) {Consulta bilingüe\\del usuario};
\node[caja,left=of query] (retr) {Recuperar\\contexto top-$k$};
\node[caja,left=of retr] (prompt) {Prompt con\\contexto y consulta};
\node[caja,below=0.9cm of retr] (slm) {Qwen2.5-3B\\cuantizado + QLoRA};
\node[caja,right=of slm] (answer) {Respuesta en\\Quechua Collao};
\draw[arrow] (docs)--(chunk);
\draw[arrow] (chunk)--(embed);
\draw[arrow] (query)--(retr);
\draw[arrow] (retr)--(prompt);
\draw[arrow] (prompt.south) |- (slm.west);
\draw[arrow] (slm)--(answer);
\draw[arrow] (embed.south)--(query.north);
\end{tikzpicture}
\par\vspace{0.12cm}
\begin{center}\scriptsize Índice documental local; generación offline.\end{center}
\fuente{Plan de tesis, arquitectura propuesta y factibilidad técnica.}
\end{frame}

% 17. Método
\begin{frame}{6 | Método: diseño de artefacto + experimento}
\begin{columns}[T,onlytextwidth]
\column{0.47\textwidth}
\begin{block}{Enfoque metodológico}
\small
\begin{itemize}
\item Investigación aplicada, cuantitativa y experimental.
\item Design Science Research (DSR): problema, objetivos, diseño, demostración, evaluación y comunicación.
\item Artefacto: pipeline local de recuperación y generación con adaptador.
\item Comparación controlada usando consultas y datos separados por partición.
\end{itemize}
\end{block}
\column{0.49\textwidth}
\begin{tcolorbox}[colback=bpLightGray,colframe=bpBlue,title={Configuraciones}]
\small
\textbf{E1} \quad SLM base cuantizado\\[1mm]
\textbf{E2} \quad SLM base + RAG local\\[1mm]
\textbf{E3} \quad SLM + RAG local + QLoRA\\[2mm]
Cada ejecución debe compartir el banco de consultas, parámetros comparables y registro de hardware.
\end{tcolorbox}
\end{columns}
\vspace{0.1cm}
\centering\scriptsize E1 establece línea base; E2 aísla el aporte de RAG; E3 permite estimar el aporte adicional de QLoRA.
\fuente{Plan de tesis, metodología y diseño experimental.}
\end{frame}

% 18. Indicadores
\begin{frame}{6 | Evaluación: qué se medirá}
\small
\begin{tabularx}{\textwidth}{>{\bfseries}p{2.6cm} X p{4.1cm}}
\toprule
Dimensión & Indicadores propuestos & Lectura\\
\midrule
Calidad lingüística & chrF++, SacreBLEU y revisión de fluidez/variedad & ¿La respuesta es comprensible y compatible con Quechua Collao?\\
Fidelidad / RAG & Faithfulness, context precision y answer relevancy (RAGAS) & ¿La respuesta sigue la evidencia y el contexto recuperado?\\
Eficiencia & Latencia, tokens/s, RAM y VRAM & ¿Puede operar bajo las restricciones offline?\\
\bottomrule
\end{tabularx}
\vspace{0.23cm}
\begin{tcolorbox}[colback=bpLightGray,colframe=bpGold]
\centering\small\textbf{Protocolo necesario:} fijar el banco de consultas, documentar versiones y parámetros, repetir mediciones y reportar limitaciones de las métricas automáticas.
\end{tcolorbox}
\fuente{Plan de tesis, métricas e instrumentos.}
\end{frame}

% 19. Factibilidad y riesgos
\begin{frame}{6 | Factibilidad, recursos y riesgos}
\begin{columns}[T,onlytextwidth]
\column{0.49\textwidth}
\begin{block}{Decisión técnica}
\small
\begin{itemize}
\item Evitar entrenamiento desde cero; partir de Qwen2.5-3B-Instruct.
\item Inferencia local cuantizada y recuperación con índice embebido.
\item Entrenar QLoRA en sesiones puntuales con GPU disponible si la máquina local no alcanza.
\item Mantener chunks y top-$k$ acotados para controlar memoria.
\end{itemize}
\end{block}
\column{0.47\textwidth}
\begin{tcolorbox}[colback=bpLightGray,colframe=bpRust,title={Riesgos y respuesta}]
\small
\textbf{OCR ruidoso} $\rightarrow$ limpieza y trazabilidad.\\[1mm]
\textbf{Mezcla de variedades} $\rightarrow$ etiquetar y filtrar dialecto.\\[1mm]
\textbf{Datos sintéticos débiles} $\rightarrow$ validar contra recursos documentales.\\[1mm]
\textbf{Memoria limitada} $\rightarrow$ cuantización, inferencia CPU y adaptación acotada.
\end{tcolorbox}
\end{columns}
\vspace{0.08cm}
\centering\scriptsize La factibilidad respalda un prototipo experimental; no demuestra de antemano su rendimiento.
\fuente{Plan de tesis, estudio de factibilidad y gestión de riesgos.}
\end{frame}

% 20. Estructura tentativa
\begin{frame}{Cómo se convierte en una tesis}
\small
\begin{columns}[T,onlytextwidth]
\column{0.47\textwidth}
\begin{enumerate}
\item \textbf{Introducción}: contexto, problema, objetivos y justificación.
\item \textbf{Marco teórico y estado del arte}: conceptos, recursos y antecedentes.
\item \textbf{Metodología y arquitectura}: diseño del artefacto y protocolo.
\end{enumerate}
\column{0.47\textwidth}
\begin{enumerate}\setcounter{enumi}{3}
\item \textbf{Implementación}: corpus, índice, SLM y adaptadores.
\item \textbf{Resultados y evaluación}: E1/E2/E3, métricas y análisis.
\item \textbf{Conclusiones y trabajo futuro}: respuesta a objetivos, límites y siguientes pasos.
\end{enumerate}
\end{columns}
\vspace{0.3cm}
\begin{tcolorbox}[colback=bpLightGray,colframe=bpGold]
\centering\small\textbf{Nota:} la organización de seis capítulos es tentativa y se ajustará al reglamento y a la evolución del proyecto.
\end{tcolorbox}
\fuente{Plan de tesis, estructura tentativa.}
\end{frame}

% 21. Puntos para cerrar antes de sustentar
\begin{frame}{Control de calidad: decisiones que faltan cerrar}
\small
\begin{enumerate}
\item \textbf{Corregir bibliografía de QueEn:} clave \texttt{Chen2024} actualmente remite a otro trabajo; confirmar autores y metadatos de arXiv:2412.05184.
\item \textbf{Alinear objetivo general e indicadores:} explicitar en el objetivo o en los específicos cómo se observarán calidad lingüística y eficiencia.
\item \textbf{Precisar evaluación lingüística:} quién revisa la variante Collao, con qué rúbrica y cómo se resolverán desacuerdos.
\item \textbf{Fijar protocolo de datos:} fuentes permitidas, particiones sin fuga, criterios de limpieza y registro de versiones.
\item \textbf{Separar plan y resultados:} hipótesis, métricas y mejoras se presentan como propuestas hasta completar experimentos.
\end{enumerate}
\vspace{0.15cm}
\begin{tcolorbox}[colback=bpLightGray,colframe=bpRust]
\centering\small Estas observaciones sirven para fortalecer el plan; no modifican el documento fuente de la tesis.
\end{tcolorbox}
\end{frame}

% 22. Referencias
\begin{frame}{Referencias principales verificadas}
\scriptsize
\setbeamertemplate{bibliography item}{}
\begin{thebibliography}{9}
\bibitem{Orcotoma2025} Orcotoma Mormontoy, G. C., Leon Nuñez, L. y Espetia Huamanga, H. (2025). \emph{Towards the Creation of a Collao Quechua--Spanish Parallel Corpus Using Optical Character Recognition}. GlobalNLP. \url{https://aclanthology.org/2025.globalnlp-1.1/}
\bibitem{Court2024} Court, S. y Elsner, M. (2024). \emph{Shortcomings of LLMs for Low-Resource Translation: Retrieval and Understanding Are Both the Problem}. WMT. \url{https://aclanthology.org/2024.wmt-1.125/}
\bibitem{Chen2024} Chen, J. et al. (2024). \emph{QueEn: A Large Language Model for Quechua-English Translation}. arXiv:2412.05184. \url{https://arxiv.org/abs/2412.05184}
\bibitem{Hevia2025} Hevia, J. S., Arredondo, F. y Kumar, V. (2025). \emph{Towards an Efficient, Customizable, and Accessible AI Tutor}. PMLR 273:250--254. \url{https://proceedings.mlr.press/v273/hevia25a.html}
\bibitem{Duran2026} Duran, M. y Silberztein, M. (2026). \emph{Semi-Automatic Construction of a Quechua-Spanish Dictionary}. LoResMT. \url{https://aclanthology.org/2026.loresmt-1.9/}
\bibitem{Walusimbi2026} Walusimbi, A. et al. (2026). \emph{Arapai: An Offline-First LLM Architecture for Adaptive Learning in Low-Connectivity Environments}. arXiv:2603.03339. \url{https://arxiv.org/abs/2603.03339}
\end{thebibliography}
\vspace{-0.1cm}
\centering\tiny Fuente de estructura y planteamiento: Plan de Tesis; proceso de referencia: presentación Seminario II adjunta.
\end{frame}

\end{document}









